% Slide 18 — Resultado principal: política única vs seleção por perfil
\begin{frame}{Política única vs. seleção por perfil}
  \centering
  \tabstyle
  \resizebox{\linewidth}{!}{%
  \begin{tabular}{@{}lrrrr@{}}
    \toprule
    \textbf{Estratégia} & \textbf{CTI médio (R\$)} & \textbf{NS médio} & \textbf{Red.\ vs.\ A1 (\%)} \\
    \midrule
    A1/A2: política única (EOQ) & 628,42 & 0,942 & 0,0 \\
    \rowcolor{accenteal!15} B: seleção por perfil & \textbf{589,62} & 0,883 & \textbf{+6,2} \\
    C$^\dagger$: oráculo por série & 347,82 & 0,700 & +44,6 \\
    \bottomrule
  \end{tabular}}
  \vspace{0.6em}

  \scriptsize $^\dagger$Oráculo: conhecimento perfeito por série -- referência
  exploratória, \textbf{não} é estratégia operacional.
  \vspace{0.55em}

  \begin{alertblock}{Leitura cautelosa}
    \small A seleção por perfil reduz CTI médio em \highlight{6,2\%}, mas com
    queda real de NS médio (0,942 $\to$ 0,883). É um \textbf{\textit{trade-off}
    custo-serviço mensurável}, não um ganho puro.
  \end{alertblock}
  \note{
    Esta é a comparação central do Experimento 2. A estratégia A é a
    política única global -- neste experimento, o EOQ é simultaneamente a
    melhor opção viável e o baseline operacional comum no varejo, então A1
    e A2 coincidem. A estratégia B seleciona a política por perfil
    operacional: na prática, isso substitui o EOQ pelo SA apenas no perfil
    Sparse High Impact, que concentra 80% das séries. O resultado é uma
    redução de 6,2% no CTI médio, mas acompanhada de queda real do nível de
    serviço, de 0,942 para 0,883 -- ainda acima do piso de viabilidade de
    0,70. Esse achado não deve ser lido como ganho puro, mas como evidência
    de que a seleção contextual produz diferença econômica mensurável, com
    uma contrapartida de serviço que precisa ser gerenciada. A estratégia C,
    o oráculo por série, é só uma referência exploratória: mostra o limite
    teórico de 44,6% de redução, mas ao custo de operar exatamente no piso
    de NS aceitável -- não é algo que se possa implementar diretamente.
  }
\end{frame}
